\documentclass[a4paper,11pt]{article}


\usepackage{amssymb,amsmath}
\usepackage[frenchb]{babel}
\usepackage[T1]{fontenc}
\usepackage[utf8]{inputenc}
\usepackage{hhline}
\usepackage{enumerate}
\usepackage{a4wide}
\usepackage{graphicx}

\newcommand{\T}{\mathbb{T}}
\newcommand{\N}{\mathbb{N}}
\newcommand{\Z}{\mathbb{Z}}
\newcommand{\Q}{\mathbb{Q}}
\newcommand{\R}{\mathbb{R}}
\newcommand{\C}{\mathbb{C}}
\newcommand{\K}{\mathbb{K}}


\def\om{\omega}
\def\eps{\varepsilon}
\def\DR{\mathcal{C}^\infty_c(\R)}

\DeclareMathOperator{\supp}{supp}

% lignes epaisses
\def\ligne#1{\leaders\hrule height #1\linethickness \hfill}

%-------------------environnements-----------------------------
\newcount\numexo
\numexo=1

\newenvironment{exercice}{\noindent\textbf{Exercice \number\numexo.---}}{\global\advance\numexo by 1}
\newenvironment{pb}{\noindent\textbf{Probl\`eme}}


\frenchspacing



%\usepackage{lmodern}
%
\usepackage{scrextend}
%%%%%%%%%%%%%%%%%%%%%%%%%%%%%%

%
\usepackage{helvet}
\renewcommand{\familydefault}{\sfdefault}
%\usepackage{rotating}
%
\thispagestyle{empty}
%
%\usepackage{inconsolata}
%
%
%\newcommand{\bigsize}{\fontsize{16pt}{20pt}\selectfont}

\usepackage{setspace}
\onehalfspacing


%%%%%%%%%%%%%%%%%%%%%%%%%%%%%%%%%%%%%%%%%%%%%%%%%%%%%%
% quelques symboles utiles en mathematiques
%%%%%%%%%%%%%%%%%%%%%%%%%%%%%%%%%%%%%%%%%%%%%%%%%%%%%%


%\newcommand{\R}{\mathbb{R}}
%\newcommand{\PP}{\mathbb{P}}
%\newcommand{\E}{\mathbb{E}}
%\newcommand{\N}{\mathbb{N}}
%\newcommand{\Q}{\mathbb{Q}}
%\newcommand{\Z}{\mathbb{Z}}

\newcommand{\sh}{{\mathop{\rm sh}}}
\newcommand{\ch}{{\mathop{\rm ch}}}
%\newcommand{\th}{{\mathop{\rm th}}}%commande deja d馭inie
\newcommand{\var}{\mathop{\rm Var}}
\newcommand{\cov}{\mathop{\rm Cov}}

                



\oddsidemargin  -0.25cm
\evensidemargin -0.25cm
\textwidth      17cm
%\headheight     0.in
%\headheight     -.25in
\topmargin      -2.5cm
\textheight     27cm





%%%%%%%%%%%%%%%%%%%%%%%%%%%%%%%%%%%%%%%%%%%%%%%%%


% les lettres grecques
\def\a{\alpha}
\def\b{\beta}
\def\l{\lambda}
%\newcommand{\a}{\alpha}\newcommand{\b}{\beta}\newcommand{\l}{\lambda}
%deja definies ailleurs
\newcommand{\g}{\gamma}
\newcommand{\m}{\mu}
\newcommand{\G}{{\Gamma}}



%%%%%%%%%%%%%% pour les lettres calligraphi仔s

\newcommand{\Arond}{{\mathcal A}}
\newcommand{\Brond}{{\mathcal B}}
\newcommand{\Crond}{{\mathcal C}}
\newcommand{\Drond}{{\mathcal D}}
\newcommand{\Erond}{{\mathcal E}}
\newcommand{\Frond}{{\mathcal F}}
\newcommand{\Grond}{{\mathcal G}}
\newcommand{\Hrond}{{\mathcal H}}
\newcommand{\Irond}{{\mathcal I}}
\newcommand{\Jrond}{{\mathcal J}}
\newcommand{\Krond}{{\mathcal K}}
\newcommand{\Lrond}{{\mathcal L}}
\newcommand{\Mrond}{{\mathcal M}}
\newcommand{\Nrond}{{\mathcal N}}
\newcommand{\Orond}{{\mathcal O}}
\newcommand{\Prond}{{\mathcal P}}
\newcommand{\Qrond}{{\mathcal Q}}
\newcommand{\Rrond}{{\mathcal R}}
\newcommand{\Srond}{{\mathcal S}}
\newcommand{\Trond}{{\mathcal T}}
\newcommand{\Urond}{{\mathcal U}}
\newcommand{\Vrond}{{\mathcal V}}
\newcommand{\Wrond}{{\mathcal W}}
\newcommand{\Xrond}{{\mathcal X}}
\newcommand{\Yrond}{{\mathcal Y}}
\newcommand{\Zrond}{{\mathcal Z}}












 \def \qed {{\hspace*{\fill}$\halmos$\medskip}}
 \def \Z {{\mathbb Z}}
 \def \R {{\mathbb R}}
 \def \C {{\mathbb C}}
 \def \N {{\mathbb N}}
 \def \P {{\mathbb P}}
 \def \E {{\mathbb E}}
 \def \Q {{\mathbb Q}}
 \def \ra {\rightarrow}
 \def \da {\downarrow}
 \def \ua {\uparrow}

 \def \cO {{\mathcal O}}
 \def \cW {{\mathcal W}}
 \def \cD {{\mathcal D}}
 \def \cL {{\mathcal L}}
 \def \cI {{\mathcal I}}
 \def \cJ {{\mathcal J}}
 \def \cR {{\mathcal R}}
 \def \cA {{\mathcal A}}
 \def \cM {{\mathcal M}}
 \def \cN {{\mathcal N}}
 \def \cE {{\mathcal E}}
 \def \cP {{\mathcal P}}
 \def \cQ {{\mathcal Q}}
 \def \cB {{\mathcal B}}
 \def \cF {{\mathcal F}}
 \def \cG {{\mathcal G}}
 \def \cC {{\mathcal C}}
 \def \cV {{\mathcal V}}
 \def \cH {{\mathcal H}}
 \def \cT {{\mathcal T}}
 \def \cZ {{\mathcal Z}}
 \def \cS {{\mathcal S}}



\begin{document}
\title{Idées de recherche}
\author{Romain Dugast}
\maketitle


\section{skein}

To go through walls in parametric families of almost complex structures, use skein valued count from tobias Ekholm.
Requires perturbation scheme -> probably close to what Russel had in mind.


\section{mcduff manifolds}

good examples of magnetic geodesic flows are given by mcduff domains (cf augustin moreno). Can we study such geodesic flow with sutured manifolds techniques ?


\section{SSMS}

We want to put a contact structure on moduli spaces. Assuming integrability of the complex
structure, this is doable. But then what is the link between the Reeb flow on the moduli space and the initial Reeb flow ?
Can we also extend this to acs homotopic to integrable ones ? How general is it ?

\section{Polysymplectic geometry}

Is there an analog of polysymplectic geometry for contact manifolds ?
Can we extend polysymplectic geometry to higher dimensional manifolds using Clifford Algebras? (Apparently yes but the compactness theory gets really hard)

\section{examples for broken books}

Classes of examples :\\
- Toric domains, convex toric domains (corresponds to manifold with a $T^2$ action,
can be represented by a polytope in the first quadrant)
- Boundary of plumbing of surfaces, fillability studied by ana reichtmann and others. 
Maybe on can construct a broken book on a plumbing ? 
However it seems like the boundary is often a lens space/a simple space.

\section{Acu}
Higer dimensional theory for open books, broken books, sutured contact manifolds. 
Transversality theory should work mostly the same, by decomposition the normal bundle. 
However, compactness may be different, with wilder possible behaviors. Could OBG still work ?
Otherwise, skein valued count, see 1. Also, can one extract topological information in that case, for sutured contact manifolds ?
THat is not so clear. Even in general, this method for extracting topological information is unclear. 

\section{examples broken book by branched cover}
to construct an example, use the construction of the counterexample and try to extract a maximal nontrivial knot from this.
Is this possible ? This does not seem very likely, but who knows ?

\section{symplectic reduction and holomorphic curves}

Can we use holomorphic curves techniques on symplectic reductions ? I think it was already done in some cases. 
Maybe I could try to use intersection theory ? 
In a simple case : take a circle action on a 6 manifold, the symplectic reduction is a 
4 dimensional manifold. Now we want it to have cylindrical ends. 
How do symplectic fillings behave with respect to symplectic reduction ? Does symplectic reduction induce contact reduction on the boundary ?
That would be cool. Is the symplectic reduction of the symplectization with a good action the symplectization of the contact reduction ?
What to do with such reduction ? What kind of result do we expect ?


\section{broken books as link singularities}

Can broken book arrise as some "link singularities" ? Maybe for real polynomials ?
Tbf that does not seem very likely as there is no symmetry. Although symmetry comes from 
the fact that the singularity is complex in the classical case, so that may not be so bad.


\section{Counterexample for the sutured conjecture}

Using work of Nathalie Wahl, we have "computable" formulas for the trace of whitehead torsion.
If we can find a sutured contact manifold, with no closed reeb orbit, and whose Dennis trace does not vanish,
that gives a counterexample to the product conjecture for sutured manifolds in higher dimension.
But this method does not seem very efficient: in order to construct a Reeb field with no closed periodic orbit,
you probably need to know the manifold... 
Another idea would be to suppose that the homotopy sutured conjecture is true but not the topological one.
Then one would need to use cut and past techniques on a product sutured contact manifold, replacing a piece by a homotopy equivalent 
one, without adding closed orbits. This seems plausible? But we need to find explicit contact structures on somehow topologically "exotic" manifolds 
such as the Whitehead space or something like that. 
In order to ensure the whole manifold is not a product, it might be interesting to use the torsion thing.
I don't know if there would be a more direct argument.

\section{Topological information from holomorphic curves}


See recent work of Zenghi Zou. Using Rabinowitz Floer homology gives much more information
Could I get Floer homological information from my foliation by holomorphic curves ? THat does not seem easy.
Also he uses a trick to get simple connectedness using the universal cover + gromov witen invariant. 
Maybe using my moduli space instead of gromov witten invariants, I could get the same results ? 


\section{Symbolic dynamics and broken book decomposition}


How to use symbolic dynamics from broken book decomposition ?
Can we use Markov partitions/shift of finite type?
Otherwise, could we use Markov chains with probability given by the area of rigid pgaes
in order to extract dynamical informations about the Reeb flows ?


\section{RFH}

Can we define Rabinowitz Floer Homology ? What would be the goal ?
Can you define Rabinowitz Floer homology for a general contact manifold in a symplectization ?
What about symplectic homology ? Does it work when there are negative ends?
 



\end{document}